\section{Related Work}
\label{sec:related}

\subsection{Appearance Matching and Local Correspondence}
Optimization-based neural style transfer separates two roles for pretrained convolutional features. A content objective constrains the representation of the carrier image, whereas a style objective matches feature correlations to a reference~\cite{gatys2016}. Feed-forward synthesis reduces the cost of producing an output by learning an image transformation, often using perceptual objectives~\cite{johnson2016}. AdaIN instead constructs a style-conditioned feature representation by matching channel-wise statistics, and WCT matches covariance through whitening and coloring~\cite{huang2017,li2017}. These mechanisms offer different ways to specify appearance, rather than interchangeable implementations of a single algorithm.

The spatial aggregation in Gram matrices and channel moments motivates explicit treatment of local correspondence. AdaAttN calculates attention-weighted local statistics and performs point-wise normalization~\cite{liu2021}. Its use of shallow and deep features also differs from a single-layer feature-matching implementation. We therefore distinguish the published methods from the local controls used in this study: the latter test whether an additional edge constraint changes an optimization result, not whether the entire proposed framework outperforms the original trained systems. This distinction is necessary because reconstruction architecture, initialization, preprocessing, and optimization budget can themselves affect the output.

\subsection{Perceptual Control and Multi-Scale Stylization}
The ability to control appearance attributes predates the present framework. Gatys et al.\ introduced controls over spatial location, color, and spatial scale, including combinations of styles from different references~\cite{gatys2017control}. LapStyle uses a learned drafting-and-revision architecture with a Laplacian pyramid to address global and local appearance at different resolutions~\cite{lapstyle2021}. These studies establish both perceptual-factor control and multi-scale processing as relevant precedents.

Our formulation uses fixed Gaussian filters to construct high-frequency and middle-frequency inputs to two losses, together with a feature-moment appearance term. It does not introduce a new pyramid representation or learn a semantic factorization. Its intended utility is an explicit set of weighted objectives whose effects can be inspected on a product-image task. Because Gaussian filters have overlapping frequency responses and visual motifs occupy multiple scales, separate coefficients alone do not guarantee independent attributes. Loss-removal images are consequently interpreted as interventions on an objective, with cross-attribute effects remaining an empirical question.

\subsection{Cultural References and Product-Image Design}
Culture-specific stylization and heritage digitization address related but distinct problems. ChipGAN provides an adversarial approach to Chinese ink-wash painting transfer~\cite{he2018}, whereas the Dunhuang project provides resources and tasks for computational heritage protection~\cite{yu2022}. The existence of such resources does not imply that their annotations, rights, or cultural interpretations transfer to an arbitrary collection of product-design references. In particular, the current five reference images form a small experimental collection rather than a representative cultural dataset.

We study reference-conditioned synthesis on ceramic tea-ware images. The product silhouette supplies the content target and a cultural-reference image supplies appearance and edge targets. This formulation differs from conservation, historical reconstruction, three-dimensional texture mapping, and evaluation by cultural specialists. The edge constraint is a measurable design prior within the image coordinate system; it does not encode the historical meaning of a lotus, deer, or porcelain motif. The practical question is whether this prior helps retain reference structure without unacceptable interference with the carrier, not whether edge alignment alone authenticates a design.

\subsection{Measurement and Culturally Situated Evaluation}
Perceptual and distributional metrics address different comparison problems. LPIPS uses deep representations calibrated against perceptual judgments~\cite{zhang2018lpips}, while FID compares distributions of image features~\cite{heusel2017fid}. A sum of uncalibrated VGG feature mean-squared errors is neither of these measures. Its scale depends on preprocessing, layer selection, and feature magnitudes. Similarly, a relative Gram-distance ratio concerns appearance alignment and does not directly measure a motif's cultural acceptability.

Culturally situated evaluation extends beyond benchmark agreement. Qadri et al.\ examine dialogue involving artists and art-world experts as a way to interpret generative tools in their creative contexts~\cite{qadri2024}. This perspective motivates a boundary in our evaluation: the CFM edge sub-score describes agreement with the same structural prior used by the optimizer, while independent cultural judgment requires separate evidence. The present computational comparisons can expose differences between diagnostics, but cannot establish the superiority of an unvalidated cultural metric or a general failure of FID and LPIPS.
